\vspace{0.5em}\noindent\textbf{A Few Takeaways from My Software Engineering Journey: From
Writing Code to Breaking the System and Fixing It
Myself}\par

\vspace{0.5em}\noindent\textbf{Introduction}\par

I once believed that the main responsibility of a Software Engineer was
to receive requirements, write code, complete features, and push the
source code to a repository. However, after working directly on a
relatively complete system, I realized that writing functional code is
only the first step.

The more difficult challenge is ensuring that the code operates reliably
within a real system containing multiple interconnected components,
including a frontend, backend APIs, databases, message queues,
containers, Kubernetes, CI/CD pipelines, and AWS infrastructure.

While developing the AI-Powered Internship Application Tracker, I worked
with FastAPI, React, PostgreSQL, Redis, DynamoDB, Amazon SQS, Docker,
Kubernetes, GitHub Actions, and several AWS services. This experience
did not simply introduce me to more technologies. It also changed how I
understand software development, deployment, failure, and recovery.

\vspace{0.5em}\noindent\textbf{Running Locally Does Not Mean Running Successfully in
Production}\par

In a local environment, system components are usually close together and
relatively easy to control. The database may run inside Docker Compose,
Redis may share the same network, and environment variables may be
stored in a local \texttt{.\allowbreak{}env} file.

When the system is moved to Kubernetes or AWS, communication becomes
more complex. A service must use the correct DNS name, namespace, port,
and network configuration to reach another service. A pod connecting to
a database requires the correct connection string, security group, and
network route. Applications in private subnets also need an appropriate
outbound path to access external services.

This experience helped me understand that an application does not
operate based on source code alone. Its reliability also depends on the
environment surrounding it.

\vspace{0.5em}\noindent\textbf{Docker and the Runtime
Environment}\par

Docker helps standardize the runtime environment across different
machines, but containerization also introduces new types of problems.

A container may be running while its health check is failing. An
application may start before its database becomes available. Environment
variables may exist on the host machine but may not be passed correctly
into the container. Port conflicts, missing volumes, Docker Desktop
settings, and WSL integration can also prevent the system from operating
even when the code has not changed.

These situations encouraged me to examine more than the application
code:

\begin{itemize}
\tightlist
\item
  Which environment variables does the container actually receive?
\item
  Are the required dependencies ready?
\item
  Are the services connected to the correct network?
\item
  Are ports exposed and mapped properly?
\item
  What behavior is the health check validating?
\end{itemize}

\vspace{0.5em}\noindent\textbf{CI/CD Is Not a Magical YAML
File}\par

GitHub Actions can automate linting, testing, image building, security
scanning, and deployment. However, a pipeline works only when all of its
assumptions are correct.

A workflow may fail because a secret is missing, a branch condition is
incorrect, an image tag is inconsistent, or the runner does not contain
the required tools. A deployment job may also be skipped because an
environment variable or execution condition has not been configured
correctly.

Through troubleshooting failed workflows, I learned that CI/CD is not
simply about creating a configuration file. Developers need to
understand what each step does, what data is passed between jobs, and
what state the destination system must be in before deployment can
proceed.

\vspace{0.5em}\noindent\textbf{IAM and the Importance of Verifying
Permissions}\par

AWS IAM was one of the areas that required the most troubleshooting.
Initially, I assumed that a missing permission could be fixed by adding
a policy containing the required action.

In practice, AWS permissions also depend on trust policies, principals,
conditions, repositories, branches, resource ARNs, and whether a
particular action supports resource-level permissions.

A policy may appear secure and carefully restricted but still fail to
operate. Some actions require the resource to be set to \texttt{"*"},
while others can be limited to a specific ARN.

Instead of continuing to guess, I began using policy simulation tools to
verify each action. The main lesson was straightforward:

\begin{quote}
Permissions should not be guessed. Each required action should be
verified to determine whether it is allowed or denied.
\end{quote}

\vspace{0.5em}\noindent\textbf{Kubernetes: Running Does Not Always Mean
Healthy}\par

Kubernetes may report a pod as \texttt{Running}, but that status only
confirms that the container currently exists. It does not guarantee that
the application inside it is functioning correctly.

A running pod may still fail to connect to Redis, wait for an unfinished
migration, return no API response, or lack outbound network access.
Therefore, system validation should include more than pod status:

\begin{itemize}
\tightlist
\item
  Application logs.
\item
  Readiness and liveness probes.
\item
  Kubernetes Services and Ingress.
\item
  Database migrations.
\item
  Network connectivity.
\item
  Dependent services.
\end{itemize}

This taught me not to rely only on infrastructure status. The actual
behavior of the application must also be tested through endpoints and
real user flows.

\vspace{0.5em}\noindent\textbf{The NAT Gateway Incident and Network
Architecture}\par

One of the most memorable incidents occurred when I disabled the NAT
Gateway to reduce operating costs. Afterward, pods in private subnets
could no longer access the Internet, some connections to AWS services
were interrupted, and deployment processes began to fail.

Because the source code, IAM configuration, and security groups had not
changed, the cause was initially difficult to identify. After reviewing
the network architecture, I discovered that the outbound route had
disappeared when the NAT Gateway was disabled.

This incident gave me a much clearer understanding of public subnets,
private subnets, Internet Gateways, NAT Gateways, route tables, and VPC
Endpoints. Concepts that had previously existed only in architecture
diagrams became practical knowledge after I experienced and resolved the
failure myself.

\vspace{0.5em}\noindent\textbf{Database Concurrency and
Idempotency}\par

Many backend problems do not appear when only one user is testing the
system. They become visible when two or more requests arrive at nearly
the same time.

For example, two requests may both check that a record does not exist
and then attempt to insert it. Without proper constraints and
transaction handling, duplicate data may be created.

To reduce these risks, the system can use:

\begin{itemize}
\tightlist
\item
  Database unique constraints.
\item
  Transactions and rollback handling.
\item
  HTTP 409 responses for conflicts.
\item
  Optimistic locking.
\item
  Idempotency keys for important operations.
\end{itemize}

An internship application should not be created twice simply because the
user clicked a button multiple times or because the client retried the
request.

\vspace{0.5em}\noindent\textbf{Message Queues and the Transactional Outbox
Pattern}\par

Another challenge appears when a system must update a database and
publish a message to a queue as part of the same logical operation.

If the database commit succeeds but message delivery fails, the data
exists while the event is lost. Sending the message first is also unsafe
because the database transaction may later fail.

The transactional outbox pattern addresses this issue by storing both
the main data and its event within the same database transaction. A
separate dispatcher then reads the outbox records and sends them to
Amazon SQS. Failed deliveries can be retried, while consumers should
support deduplication to prevent the same event from being processed
repeatedly.

This is an example of an engineering effort that users may never see
directly but that significantly improves system reliability.

\vspace{0.5em}\noindent\textbf{A Systematic Debugging
Approach}\par

The most valuable lesson was not memorizing a large number of commands.
It was learning how to divide a problem into smaller layers.

When an error occurs, I now begin with questions such as:

\begin{enumerate}
\def\labelenumi{\arabic{enumi}.}
\tightlist
\item
  Is the issue located in the application or the infrastructure?
\item
  Did it occur before or after deployment?
\item
  Are the container and pod running?
\item
  Is the service exposing the correct port?
\item
  Can DNS resolve the required hostname?
\item
  Does the network provide the necessary route?
\item
  Does IAM allow the required action?
\item
  Are the database and other dependencies ready?
\item
  What was the last step that completed successfully?
\end{enumerate}

Instead of changing several components at once, I try to identify the
first point of failure in the processing chain. This approach makes
troubleshooting more structured and reduces the time spent making random
changes.

\vspace{0.5em}\noindent\textbf{Conclusion}\par

This journey significantly changed how I understand Software
Engineering.

Previously, my first response to a system failure was to modify the
code. I now understand that a failure may originate from configuration,
environment variables, containers, networking, databases, permissions,
infrastructure, or CI/CD.

I also consider questions beyond whether a feature simply works:

\begin{itemize}
\tightlist
\item
  What happens when a request is submitted twice?
\item
  What happens if the database commits but the queue fails?
\item
  How are simultaneous updates handled?
\item
  How does the system recover when a pod restarts?
\item
  Can a failed deployment be rolled back safely?
\end{itemize}

A Software Engineer is not someone who never causes a system failure.
More importantly, a Software Engineer learns from each failure,
understands the system more deeply, resolves issues systematically, and
designs the system so that the same failure is less likely to occur
again.

\vspace{0.5em}\noindent\textbf{Publication Information}\par

\begin{itemize}
\tightlist
\item
  \textbf{Topic:} Practical lessons from developing, deploying, and
  operating a software system.
\item
  \textbf{Published date:} July 28, 2026.
\item
  \textbf{Platform:} AWS Study Groups.
\item
  \textbf{Status:} Pending.
\item
  \textbf{Public link:}
  \href{https://www.facebook.com/groups/awsstudygroupfcj/permalink/2227122848052675/\#}{Blog
  3 on AWS Study Groups}
\end{itemize}
